%%
% 摘要信息
% 本文档中前缀"c-"代表中文版字段, 前缀"e-"代表英文版字段
% 摘要内容应概括地反映出本论文的主要内容，主要说明本论文的研究目的、内容、方法、成果和结论。要突出本论文的创造性成果或新见解，不要与引言相 混淆。语言力求精练、准确，以 300—500 字为宜。
% 在摘要的下方另起一行，注明本文的关键词（3—5 个）。关键词是供检索用的主题词条，应采用能覆盖论文主要内容的通用技术词条(参照相应的技术术语 标准)。按词条的外延层次排列（外延大的排在前面）。摘要与关键词应在同一页。
% modifier: 黄俊杰(huangjj27, 349373001dc@gmail.com)
% update date: 2017-04-15
%%

\cabstract{
在学习完大数据原理相关课程后\cite{lecture01}，完成本次作业（作业-2）要求构建一个小型领域知识图谱并提交简要报告与代码。具体要求如下：

\textbf{任务要求：}
\begin{itemize}
    \item 自行选择一个小范围的主题或领域数据集（例如新闻片段、科研摘要、课程材料、产品描述等），并在报告中说明数据来源与采集方式；
    \item 使用 NLP 工具（如 spaCy）对文本数据进行预处理与信息抽取，完成实体抽取与关系抽取的流程；
    \item 使用 Neo4j 构建知识图谱，将抽取得到的实体与关系导入图数据库，并通过可视化展示关键节点与子图；
\end{itemize}

\textbf{处理流程：}
\begin{enumerate}
    \item 数据准备：使用爬虫采数据集并清洗文本数据；
    \item 实体与关系抽取：使用 spaCy 的命名实体识别（NER）识别实体类别，结合规则或简单模型抽取实体间的关系；
    \item 图数据库构建：设计实体与关系的模式（labels 与 relationship types），使用 Cypher 或官方驱动将节点与边导入 Neo4j；
    \item 可视化与分析：在 Neo4j Browser 或 Neo4j Bloom 中展示关键节点，筛选并截图作为报告结果；
\end{enumerate}

\textbf{提交说明：}
\begin{itemize}
    \item 提交包内容：代码、报告与必要的数据（或数据下载说明）；
    \item 文件命名格式：学号+姓名+作业2.zip；
    \item 提交邮箱：BDSysu2026Spring@163.com；
    \item 截止时间：布置作业之后的下一个周日 24:00；
    \item 抄袭处理：严格禁止抄袭，发现抄袭将按零分处理。
\end{itemize}

}
% 中文关键词(每个关键词之间用“，”分开,最后一个关键词不打标点符号。)
\ckeywords{知识图谱，实体抽取，关系抽取，Neo4j，可视化}
% 英文关键词(每个关键词之间用,分开, 最后一个关键词不打标点符号。)
% 下面不用写
\ekeywords{}


